\documentclass[10pt]{scrartcl}
\usepackage[mdthm,colorsec,noauthor,noasy]{evan}
\usepackage{handout}
\usepackage{figure-ink}
\newcommand{\ip}[2]{\langle #1,#2\rangle}
\DeclareMathOperator{\tr}{tr}\DeclareMathOperator{\nul}{null}\DeclareMathOperator{\range}{range}\DeclareMathOperator{\spn}{span}
\newcommand{\LL}{\mathcal L}\newcommand{\PP}{\mathcal P}\newcommand{\MM}{\mathcal M}
\title{Math 110 --- Week 12}
\subtitle{Positive operators, isometries, polar decomposition, singular values (Notes Ch.~21--22, LADR 7C--7F, HW12)}
\author{}\date{}\ihead{}

% Deterministic ink rendering; preaction textures the path without redrawing labels.
\newcommand{\PacketInkPath}[2][]{%
  \begingroup\pgfmathsetseed{1965}%
  \draw[ink line,#1,preaction={draw,arrows=-,line width=.54pt,opacity=.28,
    decorate,decoration={random steps,segment length=.65pt,amplitude=.045pt}}] #2;
  \endgroup}
\tikzset{
  >={Stealth[length=4.6pt,width=2pt,inset=1.2pt]},
  figaxis/.style={ink arrow,draw=inkslate,line width=.35pt},
  figvec/.style={ink arrow,draw=inkblue},
  figvecsub/.style={ink arrow,draw=inkslate},
  figguide/.style={ink guide,draw=inkslate},
  ink plot/.style={line width=.45pt,preaction={draw,arrows=-,line width=.54pt,
    opacity=.28,decorate,decoration={random steps,segment length=.65pt,amplitude=.045pt}}}
}

\begin{document}
\maketitle

\section{Facts used from earlier weeks}

Throughout, $V$ and $W$ are finite-dimensional inner product spaces over $\FF=\RR$ or $\CC$, with $n=\dim V$ and $m=\dim W$.

\begin{enumerate}[label=(T\arabic*),leftmargin=2.6em,itemsep=8pt]
  \ii \target{T1}\tool{Rules of the inner product (\xweek{math110}{2026-10-26-trace-inner-products}{Week~9}).} For all $u,v,w\in V$ and $\lambda\in\FF$: $\langle\lambda u+v,w\rangle=\lambda\langle u,w\rangle+\langle v,w\rangle$, $\langle u,\lambda v+w\rangle=\overline{\lambda}\langle u,v\rangle+\langle u,w\rangle$, $\langle u,v\rangle=\overline{\langle v,u\rangle}$, \textbf{positivity} $\langle v,v\rangle\ge0$ for all $v\in V$, and \textbf{definiteness} $\langle v,v\rangle=0 \iff v=0$; here $\|v\|^2=\langle v,v\rangle$. \textit{Pythagoras:} if $\langle u,w\rangle=0$, then $\|u+w\|^2=\|u\|^2+\|w\|^2$.

  \ii \target{T2}\tool{Orthonormal lists and bases (\xlink{math110}{2026-11-02-orthogonality-projections}{prop-84}{Week~10, Props.~84--87}).} For any orthonormal list $e_1,\dots,e_k$, $\|a_1e_1+\dots+a_ke_k\|^2=|a_1|^2+\dots+|a_k|^2$, and the list is linearly independent; an orthonormal list of length $\dim V$ is a basis. In an orthonormal basis $e_1,\dots,e_n$, $v=\sum_{j=1}^n\langle v,e_j\rangle e_j$. Gram--Schmidt turns any basis $v_1,\dots,v_k$ into an orthonormal basis $e_1,\dots,e_k$ with $\spn(e_1,\dots,e_j)=\spn(v_1,\dots,v_j)$ for all $j$, and any orthonormal list extends to an orthonormal basis.

  \ii \target{T3}\tool{Orthogonal complements and projections (\xlink{math110}{2026-11-02-orthogonality-projections}{prop-89}{Week~10, Props.~89--90, Def.~48}).} For a subspace $U\subseteq V$: $V=U\oplus U^\perp$ and $\dim U^\perp = n-\dim U$. The orthogonal projection $P_U$ is linear, $P_Uv\in U$, $v-P_Uv\in U^\perp$, and $P_U$ is self-adjoint (\xlink{math110}{2026-11-09-adjoints-spectral-theorem}{prob-HW11-P4}{Week~11, HW11 P4}).

  \ii \target{T4}\tool{Linear maps (\xweek{math110}{2026-09-14-linear-maps-matrices}{Week~3}).} A linear map is uniquely determined by arbitrary prescribed values on a basis. Rank--nullity: $\dim\nul T+\dim\range T=\dim V$ for $T\in\LL(V,W)$; an injective operator is invertible.

  \ii \target{T5}\tool{Adjoints (\xlink{math110}{2026-11-09-adjoints-spectral-theorem}{def-49}{Week~11, Def.~49, Prop.~93, LADR 7.5, 7.7}).} $\langle Tv,w\rangle=\langle v,T^*w\rangle$. In orthonormal bases $\MM(T^*)=\MM(T)^*$ (conjugate transpose). $(S+T)^*=S^*+T^*$, $(\lambda T)^*=\overline{\lambda}T^*$, $(T^*)^*=T$, $(TR)^*=R^*T^*$, $I^*=I$. If $T$ is invertible, $(T^*)^{-1}=(T^{-1})^*$. Also $\dim\range T^*=\dim\range T$, and two vectors having the same inner product against every basis vector are equal.

  \ii \target{T6}\tool{Self-adjoint operators (\xlink{math110}{2026-11-09-adjoints-spectral-theorem}{prop-95}{Week~11, Props.~95, 97, 101}).} Eigenvalues of a self-adjoint operator are real. If $T=T^*$ and $\langle Tv,v\rangle=0$ for all $v$, then $T=0$. If $T=T^*$ and $U$ is $T$-invariant, $T|_U$ is self-adjoint.

  \ii \target{T7}\tool{Spectral theorems (\xlink{math110}{2026-11-09-adjoints-spectral-theorem}{prop-102}{Week~11, Props.~99, 102, 104}).} Every self-adjoint operator (over $\RR$ or $\CC$) is orthogonally diagonalizable; conversely, any operator with an orthonormal eigenbasis and real eigenvalues is self-adjoint (Prop.~102(b)). Over $\CC$, $T$ is normal ($TT^*=T^*T$) $\iff T$ is orthogonally diagonalizable; if $Tv=\lambda v$ for normal $T$, then $T^*v=\overline{\lambda}v$. A square matrix $Q$ with orthonormal columns satisfies $Q^*Q=I$ and $Q^{-1}=Q^*$ (called \textbf{orthogonal} over $\RR$, \textbf{unitary} over $\CC$). Self-adjoint matrices factor as $QDQ^*$ with $D$ real diagonal.

  \ii \target{T8}\tool{Diagonal matrices (\xweek{math110}{2026-10-05-eigenvalues-pagerank}{Week~6}).} If $Te_k=d_k e_k$ on a basis $e_1,\dots,e_n$, the eigenvalues of $T$ are $\{d_1,\dots,d_n\}$ (counting multiplicity); if every $d_k=\mu$, then $T=\mu I$.

  \ii \target{T9}\tool{$2\times2$ matrices (\xweek{math110}{2026-10-05-eigenvalues-pagerank}{Week~6}).} $\det\begin{psmallmatrix}a&b\\c&d\end{psmallmatrix}=ad-bc\neq0\iff$ invertible; characteristic polynomial $x^2-(\tr A)x+\det A$.

  \ii \target{T10}\tool{Norms (\xlink{math110}{2026-10-26-trace-inner-products}{def-44}{Week~9, Def.~44, Prop.~83}).} A norm satisfies positivity $\|v\|\ge0$, definiteness $\|v\|=0\iff v=0$, homogeneity $\|\lambda v\|=|\lambda|\,\|v\|$, and the triangle inequality $\|u+v\|\le\|u\|+\|v\|$. Inner-product norms ($\|v\|^2=\langle v,v\rangle$) satisfy the parallelogram law $\|u+v\|^2+\|u-v\|^2=2\|u\|^2+2\|v\|^2$.

  \ii \target{T11}\tool{Trace (\xlink{math110}{2026-10-26-trace-inner-products}{prop-78}{Week~9, Prop.~78}).} $\tr(AB)=\tr(BA)$, so $\tr(Q^{-1}AQ)=\tr A$ is basis-independent.
\end{enumerate}

\section{This week}

Just as complex numbers factor into $z = e^{i\theta}\cdot r$ (a unit-modulus rotation times a nonnegative real stretch $r=|z|=\sqrt{\overline{z}z}$), we decompose general linear maps into an \textbf{isometry} (length-preserving) times a \textbf{positive operator} (pure orthogonal stretching).

\subsection{Positive operators (Ch.~21)}

\begin{qbox}
  For the two symmetric matrices below, does $\langle Av,v\rangle\ge0$ hold for every $v\in\RR^2$? Does either matrix admit a symmetric square root $R$ with $R^2=A$?
  \[
    A_1 = \begin{bmatrix}2&-1\\-1&2\end{bmatrix}, \qquad
    A_2 = \begin{bmatrix}1&3\\3&1\end{bmatrix}.
  \]
\end{qbox}

Among self-adjoint operators (the ``real numbers'' of $\LL(V)$), positive operators play the role of nonnegative real numbers $r\ge0$.

\begin{definition*}[Def.~52 / LADR 7C: Positive and positive definite operators]\target{def-52}
  A self-adjoint operator $T\in\LL(V)$ is \textbf{positive} (also \textbf{positive semidefinite}) if
  \[
    \langle Tv,v\rangle \ge 0 \qquad \text{for all } v\in V.
  \]
  It is \textbf{positive definite} if $\langle Tv,v\rangle > 0$ for all $v\neq0$. (Note that Axler LADR requires self-adjointness explicitly: over $\RR$, skew operators can satisfy $\langle Tv,v\rangle=0$.) A square matrix is positive (definite) if the operator $x\mapsto Ax$ on $\FF^n$ with the standard inner product is.
\end{definition*}

\begin{example*}[Ex.~75: Orthogonal projections are positive]\target{ex-75}
  Let $P_U$ be the orthogonal projection onto a subspace $U\subseteq V$, and write $v=P_Uv+(v-P_Uv)$ with $P_Uv\in U$ and $v-P_Uv\in U^\perp$ \ilink{T3}{(T3)}. Because these pieces are orthogonal,
  \[
    \langle P_U v,v\rangle = \langle P_Uv,P_Uv\rangle + \langle P_Uv,v-P_Uv\rangle = \|P_Uv\|^2 \ge 0.
  \]
  Since $P_U$ is self-adjoint \ilink{T3}{(T3)}, it is positive. It is positive definite only when $U=V$ (a nonzero $v\in U^\perp$ has $\langle P_Uv,v\rangle=0$).
\end{example*}

\begin{theorem*}[Prop.~105 / LADR 7.38: Six descriptions of a positive operator]\target{prop-105}
  For $T\in\LL(V)$, the following are equivalent:
  \begin{enumerate}
    \ii[(i)] $T$ is positive;
    \ii[(ii)] $T$ is self-adjoint and every eigenvalue of $T$ is nonnegative;
    \ii[(iii)] $T=R^2$ for some self-adjoint $R\in\LL(V)$;
    \ii[(iv)] $T=R^*R$ for some $R\in\LL(V)$;
    \ii[(v)] in some orthonormal basis, $\MM(T)$ is diagonal with all diagonal entries $\ge0$;
    \ii[(vi)] $T$ has a positive square root: $T=R^2$ for some positive $R\in\LL(V)$.
  \end{enumerate}
\end{theorem*}
\begin{ideabox}
  On an eigenvector $v$, $\langle Tv,v\rangle = \lambda\|v\|^2$, so positivity forces $\lambda\ge0$. Conversely, by the spectral theorem \ilink{T7}{(T7)}, a self-adjoint $T$ with eigenvalues $\lambda_k\ge0$ acts on an orthonormal basis by multiplying each basis vector by $\lambda_k\ge0$; replacing each $\lambda_k$ by $\sqrt{\lambda_k}\ge0$ constructs a positive square root $R$. Moreover, any $R^*R$ satisfies $\langle R^*Rv,v\rangle = \|Rv\|^2 \ge 0$.
\end{ideabox}
\begin{proof}
  \begin{steps}
    \ii \textbf{(i)$\Rightarrow$(ii).} $T$ is self-adjoint by Def.~52. If $Tv=\lambda v$ ($v\neq0$), then $\lambda\in\RR$ \ilink{T6}{(T6)} and $\lambda\|v\|^2=\langle Tv,v\rangle\ge0$, forcing $\lambda\ge0$.
    \ii \textbf{(ii)$\Rightarrow$(iii) \& (vi).} By \ilink{T7}{(T7)}, choose an orthonormal basis $e_1,\dots,e_n$ with $Te_k=\lambda_k e_k$ ($\lambda_k\ge0$). Define $R\in\LL(V)$ by $Re_k=\sqrt{\lambda_k}e_k$ \ilink{T4}{(T4)}. Since $R$ has real eigenvalues $\sqrt{\lambda_k}\ge0$ in an orthonormal basis, $R$ is self-adjoint \ilink{T7}{(T7)} with eigenvalues $\sqrt{\lambda_k}\ge0$ \ilink{T8}{(T8)}, and $R^2e_k=\lambda_k e_k=Te_k$ forces $R^2=T$. Moreover $\langle Rv,v\rangle = \sum \sqrt{\lambda_k}|\langle v,e_k\rangle|^2 \ge 0$, so $R$ is positive.
    \ii \textbf{(vi)$\Rightarrow$(iii).} A positive operator $R$ is self-adjoint by Def.~52.
    \ii \textbf{(iii)$\Rightarrow$(iv).} If $R^*=R$, then $T=R^2=R^*R$.
    \ii \textbf{(iv)$\Rightarrow$(i).} By \ilink{T5}{(T5)}, $(R^*R)^*=R^*(R^*)^*=R^*R$, so $T$ is self-adjoint; and $\langle Tv,v\rangle=\langle R^*Rv,v\rangle=\|Rv\|^2\ge0$.
    \ii \textbf{(ii)$\Leftrightarrow$(v).} In step~(ii)$\Rightarrow$(iii), $\MM(T)=\operatorname{diag}(\lambda_1,\dots,\lambda_n)$ with $\lambda_k\ge0$. Conversely, an orthonormal basis where $Te_k=d_ke_k$ ($d_k\ge0$) makes $T$ self-adjoint \ilink{T7}{(T7)} with eigenvalues $d_1,\dots,d_n\ge0$ \ilink{T8}{(T8)}.
  \end{steps}
\end{proof}

\begin{proposition*}[Prop.~106: Positive definite operators]\target{prop-106}
  For $T\in\LL(V)$, the following are equivalent:
  \begin{enumerate}
    \ii[(i)] $T$ is positive definite;
    \ii[(ii)] $T$ is self-adjoint and every eigenvalue of $T$ is strictly positive ($\lambda>0$);
    \ii[(iii)] $T=R^*R$ for some invertible $R\in\LL(V)$.
  \end{enumerate}
\end{proposition*}
\begin{proof}
  \begin{steps}
    \ii \textbf{(i)$\Rightarrow$(ii).} If $Tv=\lambda v$ ($v\neq0$), then $\lambda\|v\|^2=\langle Tv,v\rangle>0$, so $\lambda>0$.
    \ii \textbf{(ii)$\Rightarrow$(iii).} Build $Re_k=\sqrt{\lambda_k}e_k$ as in Prop.~105. Since $\sqrt{\lambda_k}>0$ for all $k$, $\nul R=\{0\}$, so $R$ is invertible \ilink{T4}{(T4)} and $T=R^2=R^*R$.
    \ii \textbf{(iii)$\Rightarrow$(i).} $R^*R$ is self-adjoint and for $v\neq0$, injectivity of $R$ forces $Rv\neq0$, hence $\langle Tv,v\rangle=\|Rv\|^2>0$.
  \end{steps}
\end{proof}

\begin{proposition*}[Algebra of positive operators]
  If $S,T\in\LL(V)$ are positive, then:
  \begin{enumerate}[(a)]
    \ii $S+T$ is positive;
    \ii $cT$ is positive for every real $c\ge0$;
    \ii $X^*TX$ is positive for every $X\in\LL(V)$;
    \ii if $T$ is invertible, then $T$ is positive definite.
  \end{enumerate}
\end{proposition*}
\begin{proof}
  (a), (b) Self-adjointness follows from \ilink{T5}{(T5)}, and $\langle(S+T)v,v\rangle=\langle Sv,v\rangle+\langle Tv,v\rangle\ge0$, $\langle cTv,v\rangle=c\langle Tv,v\rangle\ge0$.
  (c) $(X^*TX)^*=X^*T^*X=X^*TX$, and $\langle X^*TXv,v\rangle=\langle T(Xv),Xv\rangle\ge0$.
  (d) Eigenvalues are $\ge0$ (Prop.~105), and $0$ cannot be an eigenvalue since $T$ is invertible; hence all eigenvalues are $>0$ (Prop.~106).
\end{proof}

\begin{proposition*}[Quadratics in a self-adjoint operator, LADR 7C]
  If $T\in\LL(V)$ is self-adjoint and $b,c\in\RR$ satisfy $b^2<4c$, then $T^2+bT+cI$ is positive definite (and hence invertible).
\end{proposition*}
\begin{proof}
  $Q=T^2+bT+cI$ is self-adjoint by \ilink{T5}{(T5)}. In an orthonormal eigenbasis $e_1,\dots,e_n$ of $T$ with real eigenvalues $\lambda_k$, $Qe_k = \mu_k e_k$ where $\mu_k = (\lambda_k+b/2)^2 + (c-b^2/4) > 0$, so $Q$ is positive definite by Prop.~106.
\end{proof}

\begin{theorem*}[Prop.~107 / LADR 7.39, 7.40: Uniqueness of the positive square root]\target{prop-107}
  Every positive operator $T\in\LL(V)$ has a \textbf{unique} positive square root, denoted $\sqrt{T}$. On each eigenspace $E(\lambda,T)$, $\sqrt{T}$ acts as $\sqrt{\lambda}\,I$.
\end{theorem*}
\begin{ideabox}
  Any positive square root $R$ ($R^2=T$) commutes with $T=R^2$: $RT=R^3=TR$. Therefore $R$ leaves every eigenspace $E(\lambda,T)$ invariant. Inside $E(\lambda,T)$, $R^2=\lambda I$, so every eigenvalue $\mu$ of the positive operator $R|_{E(\lambda,T)}$ satisfies $\mu^2=\lambda$ and $\mu\ge0$, forcing $\mu=\sqrt{\lambda}$. Diagonalizing this self-adjoint restriction yields $R=\sqrt{\lambda}\,I$ on every $E(\lambda,T)$.
\end{ideabox}
\begin{proof}
  \begin{steps}
    \ii Existence was proved in Prop.~105(vi). Suppose $R\in\LL(V)$ is any positive operator with $R^2=T$.
    \ii If $Tv=\lambda v$, then $T(Rv)=R^3v=R(Tv)=\lambda Rv$, so $E(\lambda,T)$ is $R$-invariant. Put $R_\lambda=R|_{E(\lambda,T)}$, which is self-adjoint \ilink{T6}{(T6)} and positive ($\langle R_\lambda v,v\rangle=\langle Rv,v\rangle\ge0$).
    \ii Since $R_\lambda^2=\lambda I$, any eigenvalue $\mu$ of $R_\lambda$ satisfies $\mu^2=\lambda$ and $\mu\ge0$ (Prop.~105), forcing $\mu=\sqrt{\lambda}$. Because $R_\lambda$ is self-adjoint, $E(\lambda,T)$ has a basis of eigenvectors of $R_\lambda$ \ilink{T7}{(T7)}, all with eigenvalue $\sqrt{\lambda}$, so $R_\lambda=\sqrt{\lambda}\,I$ \ilink{T8}{(T8)}.
    \ii Since $V=\bigoplus_\lambda E(\lambda,T)$ \ilink{T7}{(T7)}, the values of $R$ on an eigenbasis are uniquely forced, so $R$ is unique \ilink{T4}{(T4)}.
  \end{steps}
\end{proof}

\begin{example*}[Ex.~76: Without positivity, square roots are not unique]\target{ex-76}
  For any $a,b\in\CC$ with $b\neq0$, the matrix $R=\begin{bmatrix}a&b\\(1-a^2)/b&-a\end{bmatrix}$ satisfies $R^2=I$. Thus $I$ on $\CC^2$ has infinitely many square roots (this entire family, $\pm I$, $\operatorname{diag}(1,-1)$, etc.), yet by Prop.~107, $I$ itself is the \emph{only} positive square root.
\end{example*}

\noindent\textbf{Computing $\sqrt{T}$ and $2\times2$ shortcut.} If $A=QDQ^*$ is an orthogonal/unitary diagonalization \ilink{T7}{(T7)} of a positive matrix $A$, then $\sqrt{A}=Q\sqrt{D}Q^*$ where $\sqrt{D}$ takes entrywise square roots. For $2\times2$ matrices, no eigenvectors are needed.

\begin{proposition*}[Square-root formula for $2\times2$ positive matrices]\target{name-square-root-2x2}
  Let $A$ be a nonzero positive $2\times2$ matrix. Then
  \[
    \sqrt{A} = \frac{A + \sqrt{\det A}\,I}{\sqrt{\tr A + 2\sqrt{\det A}}}.
  \]
\end{proposition*}
\begin{proof}
  Let $\lambda_1,\lambda_2\ge0$ be the eigenvalues of $A$ with orthonormal eigenvectors $u_1,u_2$, and set $s_k=\sqrt{\lambda_k}\ge0$ ($s_1+s_2>0$ since $A\neq0$). For $k=1,2$,
  \[
    (A+s_1s_2 I)u_k = (s_k^2+s_1s_2)u_k = s_k(s_1+s_2)u_k = (s_1+s_2)\sqrt{A}\,u_k,
  \]
  so $A+s_1s_2 I = (s_1+s_2)\sqrt{A}$ \ilink{T4}{(T4)}. By \ilink{T9}{(T9)}, $s_1s_2=\sqrt{\lambda_1\lambda_2}=\sqrt{\det A}$ and $(s_1+s_2)^2=\lambda_1+\lambda_2+2s_1s_2=\tr A+2\sqrt{\det A}$.
\end{proof}

\begin{fig}
\begin{tikzpicture}[scale=1.0]
  % Panel (a): Positive definite level set (ellipse)
  \begin{scope}[xshift=2.8cm,yshift=1.95cm,scale=1.24]
  \InkRegion[inkblue]{(0,0) [rotate=45] ellipse (1 and 0.57735)}
    \PacketInkPath[figaxis]{(-1.42,0) -- (1.55,0) node[right=2pt,font=\scriptsize] {\InkLabel{$x$}}}
    \PacketInkPath[figaxis]{(0,-1.30) -- (0,1.42) node[above=1pt,font=\scriptsize] {\InkLabel{$y$}}}
    % Orthonormal eigenvector axes (dotted reference)
    \draw[ink guide,draw=black!45] (-1.00,-1.00) -- (1.08,1.08)
      node[above right=0pt,font=\scriptsize,text=black!75] {\InkLabel{$u$}};
    \draw[ink guide,draw=black!45] (0.90,-0.90) -- (-0.95,0.95)
      node[above left=0pt,font=\scriptsize,text=black!75] {\InkLabel{$w$}};

    % Soft neutral gray wash + crisp black ellipse
      \PacketInkPath[black,line width=0.55pt,rotate=45]{(0,0) ellipse (1 and 0.57735)}

    % Semi-axis lengths
    \PacketInkPath[->,black!80,line width=0.55pt,shorten <=1.0pt,shorten >=0.8pt]{(0,0) -- (0.707,0.707)}
    \node[font=\scriptsize,text=black!85] at (0.50,0.22) {\InkLabel{$1$}};

    \PacketInkPath[->,black!80,line width=0.55pt,shorten <=1.0pt,shorten >=0.8pt]{(0,0) -- (-0.408,0.408)}
    \node[font=\scriptsize,text=black!85] at (-0.78,0.64) {\InkLabel{$1/\sqrt{3}$}};

    \fill[black] (0,0) circle (1.2pt);
    \node[font=\footnotesize,anchor=north] at (0,-1.40) {\InkLabel{(a)}};
  \end{scope}

  % Panel (b): Indefinite level set (hyperbola)
  \begin{scope}[xshift=9.6cm,yshift=1.95cm,scale=1.24]
    \PacketInkPath[figaxis]{(-1.42,0) -- (1.55,0) node[right=2pt,font=\scriptsize] {\InkLabel{$x$}}}
    \PacketInkPath[figaxis]{(0,-1.30) -- (0,1.42) node[above=1pt,font=\scriptsize] {\InkLabel{$y$}}}
    % Orthonormal eigenvector axes (dotted reference)
    \draw[ink guide,draw=black!45] (-1.00,-1.00) -- (1.08,1.08)
      node[above right=0pt,font=\scriptsize,text=black!75] {\InkLabel{$u$}};
    \draw[ink guide,draw=black!45] (-0.95,0.95) -- (0.95,-0.95)
      node[below right=0pt,font=\scriptsize,text=black!75] {\InkLabel{$w$}};

    % Dotted asymptotes w = +/- sqrt(2)*u of 4u^2 - 2w^2 = 0 (rotated 45 deg)
    \begin{scope}
      \clip (-1.28,-1.20) rectangle (1.28,1.25);
      \draw[ink guide,draw=inkslate!65,rotate=45] (-0.78,-1.103) -- (0.78,1.103);
      \draw[ink guide,draw=inkslate!65,rotate=45] (-0.78,1.103) -- (0.78,-1.103);
    \end{scope}

    % Crisp pure black hyperbola branches
    \draw[ink line,black,line width=0.55pt,smooth,domain=-1.00:1.00,samples=70,variable=\t,rotate=45] plot ({0.5*cosh(\t)},{0.7071*sinh(\t)});
    \draw[ink line,black,line width=0.55pt,smooth,domain=-1.00:1.00,samples=70,variable=\t,rotate=45] plot ({-0.5*cosh(\t)},{0.7071*sinh(\t)});

    \node[font=\scriptsize] at (1.18,0.52) {\InkLabel[inkblue]{$\langle A_2v,v\rangle>0$}};
    \node[font=\scriptsize] at (-1.12,0.48) {\InkLabel[inkred]{$\langle A_2v,v\rangle<0$}};
    \fill[black] (0,0) circle (1.2pt);
    \node[font=\footnotesize,anchor=north] at (0,-1.40) {\InkLabel{(b)}};
  \end{scope}
\end{tikzpicture}
\figcap{Level sets $\langle Av,v\rangle=1$ along the orthonormal eigenvector axes $u,w$ (dotted). (a) Positive definite $A_1$ ($\lambda=1,3$) gives the closed ellipse $u^2+3w^2=1$ with semi-axes $1$ and $1/\sqrt{3}$; (b) indefinite $A_2$ ($\lambda=4,-2$) gives the hyperbola $4u^2-2w^2=1$.}
\end{fig}

\noindent\textbf{The question answered.} $A_1$ has $\tr A_1=4$, $\det A_1=3$, eigenvalues $3,1>0$, so $A_1$ is positive definite (Prop.~106), and $\sqrt{A_1}=\frac{A_1+\sqrt{3}I}{\sqrt{4+2\sqrt{3}}}=\frac{1}{2}\begin{bmatrix}\sqrt{3}+1&1-\sqrt{3}\\1-\sqrt{3}&\sqrt{3}+1\end{bmatrix}$. Meanwhile $\det A_2=-8<0$ (eigenvalues $4,-2$), so $A_2$ is \emph{not} positive ($\langle A_2(1,-1),(1,-1)\rangle=-4$) and has no symmetric square root.

\begin{theorem*}[Prop.~108: Positive definite matrices classify all inner products on $\FF^n$]\target{prop-108}
  For every inner product on $\FF^n$, there exists a unique Hermitian positive definite matrix $G$ such that
  \[
    \langle v,w\rangle = v^{\mathsf T}G\overline{w} \qquad \text{for all } v,w\in\FF^n,
  \]
  namely $G_{ij}=\langle e_i,e_j\rangle$. Conversely, every Hermitian positive definite $G$ defines an inner product by this formula. (Over $\RR$: $\langle v,w\rangle=v^{\mathsf T}Gw$ with $G$ real symmetric positive definite.)
\end{theorem*}
\begin{proof}
  \begin{steps}
    \ii Expanding $v=\sum_i v_i e_i$ and $w=\sum_j w_j e_j$ via \ilink{T1}{(T1)} gives $\langle v,w\rangle = \sum_{i,j} v_i \overline{w_j}\langle e_i,e_j\rangle = v^{\mathsf T}G\overline{w}$ with $G_{ij}=\langle e_i,e_j\rangle$ forced by taking $v=e_i,w=e_j$.
    \ii $G_{ji}=\langle e_j,e_i\rangle=\overline{\langle e_i,e_j\rangle}=\overline{G_{ij}}$, so $G^*=G$. Under the standard inner product $\langle\cdot,\cdot\rangle_0$, for $y\neq0$ setting $v=\overline{y}\neq0$ yields $\langle Gy,y\rangle_0 = \sum_{i,j} \overline{v_j} G_{ij} v_i = v^{\mathsf T}G\overline{v} = \langle v,v\rangle > 0$.
    \ii Conversely, any Hermitian positive definite $G$ makes $\langle v,w\rangle=v^{\mathsf T}G\overline{w}$ linear in $v$. Because $G^*=\overline{G}^{\mathsf T}=G$ and a scalar equals its transpose, $\overline{w^{\mathsf T}G\overline{v}}=\overline{w}^{\mathsf T}\overline{G}v = (\overline{w}^{\mathsf T}\overline{G}v)^{\mathsf T} = v^{\mathsf T}G^*\overline{w}=v^{\mathsf T}G\overline{w}$, giving conjugate symmetry. Finally, by Step~2, $\langle v,v\rangle=\langle G\overline{v},\overline{v}\rangle_0 > 0$ for all $v\neq0$.
  \end{steps}
\end{proof}

\noindent\textbf{Where this is used: covariance matrices \& PCA (Notes Ch.~21).} For a random vector $X\in\RR^n$ with mean $\mu=\mathbb{E}X$, its covariance matrix $C=\mathbb{E}[(X-\mu)(X-\mu)^{\mathsf T}]$ is symmetric and positive because $v^{\mathsf T}Cv = \mathbb{E}[(v^{\mathsf T}(X-\mu))^2]\ge0$. Maximizing variance $v^{\mathsf T}Cv$ over unit vectors $v$ yields the top eigenvector of $C$; orthogonal eigenvectors give the \textbf{principal directions} of Principal Component Analysis (PCA).

\begin{shortcuts}
  \ii \textbf{$2\times2$ positivity:} A Hermitian $2\times2$ matrix is positive $\iff \det A\ge0$ and $\tr A\ge0$ (strict $>$ for positive definite).
  \ii \textbf{Diagonal entries:} Any positive matrix has $a_{ii}=\langle Ae_i,e_i\rangle\ge0$.
  \ii \textbf{$R^*R$ is always positive:} So are $RR^*$, $A^{\mathsf T}A$, and covariance matrices (Prop.~105).
  \ii \textbf{Zero value of a positive operator (LADR 7C):} If $T$ is positive and $\langle Tv,v\rangle=0$, then $\|\sqrt{T}v\|^2=0 \implies \sqrt{T}v=0 \implies Tv=0$.
\end{shortcuts}

\begin{trapbox}
  \ii \textbf{Positive entries $\neq$ positive matrix:} $A_2=\begin{psmallmatrix}1&3\\3&1\end{psmallmatrix}$ has all entries positive yet $\det A_2=-8<0$; $A_1$ has negative off-diagonal entries yet is positive definite.
  \ii \textbf{Entry $a_{11}\ge0$ and $\det A\ge0$ is not enough:} $\operatorname{diag}(0,-1)$ has $a_{11}=0$ and $\det=0$ yet $\tr=-1<0$. Always check $\tr A\ge0$.
  \ii \textbf{$\sqrt{A}$ is not entrywise:} taking entrywise roots gives the wrong matrix:
  \[
    \sqrt{\begin{bmatrix}5&4\\4&5\end{bmatrix}} = \begin{bmatrix}2&1\\1&2\end{bmatrix} \neq \begin{bmatrix}\sqrt{5}&2\\2&\sqrt{5}\end{bmatrix}.
  \]
  \ii \textbf{Must check self-adjointness:} $90^\circ$ rotation on $\RR^2$ has $\langle Tv,v\rangle=0\ge0$ everywhere, yet is not positive.
\end{trapbox}

\subsection{Isometries (Ch.~21)}

\begin{qbox}
  The map $S\colon\RR^2\to\RR^3$ given by $S(x,y)=(x,y,0)$ preserves all lengths: $\|S(x,y)\|=\|(x,y)\|$. Does $S^*$ preserve lengths? If $S\colon\RR^n\to\RR^n$ is an operator on the \emph{same} space that preserves lengths, must $S^*S=SS^*=I$?
\end{qbox}

\begin{definition*}[Def.~53 / LADR 7.44: Isometry]\target{def-53}
  A linear map $S\in\LL(V,W)$ is an \textbf{isometry} if $\|Sv\|=\|v\|$ for every $v\in V$. An isometric operator $S\in\LL(V)$ is called \textbf{orthogonal} over $\RR$ and \textbf{unitary} over $\CC$ (LADR 4th ed.\ 7.44 defines isometry for $S\in\LL(V,W)$ and calls isometric operators unitary in LADR 7.51).
\end{definition*}

\begin{theorem*}[Prop.~109 / LADR 7.49: Isometry $\iff S^*S=I$]\target{prop-109}
  $S\in\LL(V,W)$ is an isometry if and only if $S^*S=I_V$.
\end{theorem*}
\begin{proof}
  For all $v\in V$, $\langle(S^*S-I)v,v\rangle = \langle Sv,Sv\rangle - \langle v,v\rangle = \|Sv\|^2 - \|v\|^2$. Since $S^*S-I$ is self-adjoint \ilink{T5}{(T5)}, $\|Sv\|^2-\|v\|^2=0$ for all $v \iff S^*S-I=0$ \ilink{T6}{(T6)} $\iff S^*S=I$.
\end{proof}

\begin{theorem*}[Prop.~110 / LADR 7.49: Isometries and orthonormal lists]\target{prop-110}
  For $S\in\LL(V,W)$, the following are equivalent:
  \begin{enumerate}
    \ii[(i)] $S$ is an isometry;
    \ii[(ii)] $\langle Su,Sv\rangle = \langle u,v\rangle$ for all $u,v\in V$ ($S$ preserves angles as well as lengths);
    \ii[(iii)] $S$ maps every orthonormal basis $e_1,\dots,e_n$ of $V$ to an orthonormal list $Se_1,\dots,Se_n$ in $W$;
    \ii[(iv)] $S$ maps some orthonormal basis $e_1,\dots,e_n$ of $V$ to an orthonormal list in $W$.
  \end{enumerate}
\end{theorem*}
\begin{proof}
  \begin{steps}
    \ii \textbf{(i)$\Rightarrow$(ii).} $\langle Su,Sv\rangle=\langle S^*Su,v\rangle=\langle u,v\rangle$ by Prop.~109.
    \ii \textbf{(ii)$\Rightarrow$(iii).} $\langle Se_i,Se_j\rangle=\langle e_i,e_j\rangle=\delta_{ij}$.
    \ii \textbf{(iii)$\Rightarrow$(iv).} $V$ has an orthonormal basis \ilink{T2}{(T2)}.
    \ii \textbf{(iv)$\Rightarrow$(i).} Writing $v=\sum_{k=1}^n \lambda_k e_k$, \ilink{T2}{(T2)} applied to the orthonormal list $Se_1,\dots,Se_n$ in $W$ gives $\|Sv\|^2 = \|\sum \lambda_k Se_k\|^2 = \sum |\lambda_k|^2 = \|v\|^2$.
  \end{steps}
\end{proof}

\begin{proposition*}[Isometry inverse rule]
  Let $S\in\LL(V)$ be an isometric \textbf{operator} ($W=V$). Then $S$ is invertible with $S^{-1}=S^*$, $SS^*=I$, $S$ is normal, and $S^*$ is also an isometry. In particular, if a square matrix has orthonormal columns, then its rows are also orthonormal.
\end{proposition*}
\begin{proof}
  If $Sv=0$ then $\|v\|=\|Sv\|=0$, so $S$ is injective, hence invertible \ilink{T4}{(T4)}. From $S^*S=I$ (Prop.~109), right-multiplying by $S^{-1}$ gives $S^*=S^{-1}$, hence $SS^*=I=S^*S$ (normal) and $(S^*)^*S^*=SS^*=I$ ($S^*$ is an isometry). For a square matrix $S$, its columns are $S\delta_j$ (orthonormal), so $S^*$ is an isometry whose columns $S^*\delta_k$ (the conjugates of the rows of $S$) are orthonormal, hence the rows of $S$ are orthonormal.
\end{proof}

\begin{theorem*}[Prop.~111 / LADR 7.55: Complex isometries]\target{prop-111}
  Let $V$ be complex and $S\in\LL(V)$. Then $S$ is an isometry if and only if $S$ is normal and every eigenvalue of $S$ has absolute value $1$.
\end{theorem*}
\begin{proof}
  ($\Rightarrow$) An isometric operator is normal by the inverse rule, and if $Sv=\lambda v$ ($v\neq0$), then $\|v\|=\|Sv\|=|\lambda|\,\|v\|\implies|\lambda|=1$.
  ($\Leftarrow$) By \ilink{T7}{(T7)}, let $e_1,\dots,e_n$ be an orthonormal eigenbasis with $Se_k=\lambda_k e_k$, $|\lambda_k|=1$. Since $S$ is normal, $S^*e_k=\overline{\lambda_k}e_k$, so $S^*Se_k=|\lambda_k|^2 e_k = e_k$ for all $k$, giving $S^*S=I$ \ilink{T4}{(T4)}.
\end{proof}

\begin{example*}[Householder reflection $H=I-2uu^{\mathsf T}$]
  Let $u\in\RR^n$ be a unit vector ($\|u\|=1$, $u^{\mathsf T}u=1$) and set $H=I-2uu^{\mathsf T}$. Because $(uu^{\mathsf T})^{\mathsf T}=uu^{\mathsf T}$, we have $H^{\mathsf T}=H$ and
  \[
    H^{\mathsf T}H = H^2 = I - 4uu^{\mathsf T} + 4u(u^{\mathsf T}u)u^{\mathsf T} = I,
  \]
  so $H$ is both self-adjoint and orthogonal (Prop.~109). Geometrically, decomposing any $v\in\RR^n$ into its component $(u^{\mathsf T}v)u$ along $u$ and $w=v-(u^{\mathsf T}v)u\in u^\perp$, we have $Hu=-u$ and $Hw=w$: $H$ reflects every vector across the hyperplane $u^\perp$, with eigenvalues $+1$ (multiplicity $n-1$ on $u^\perp$) and $-1$ (multiplicity $1$ on $\spn(u)$).
\end{example*}

\begin{fig}
\begin{tikzpicture}[scale=1.05]
  % Orthogonal frame rotated by 16 degrees
  \begin{scope}[rotate=16]
    % Hyperplane u^\perp and normal line span(u)
    \PacketInkPath[black!85,line width=0.55pt]{(-2.05,0) -- (2.95,0)}
    \draw[ink guide,draw=black!45] (0,-1.40) -- (0,1.70);

    % Delicate right-angle marks at origin and projection foot W
    \PacketInkPath[black!55,line width=0.38pt]{(0.20,0) -- (0.20,0.20) -- (0,0.20)}
    \PacketInkPath[black!55,line width=0.38pt]{(2.20,0) -- (2.20,0.20) -- (2.40,0.20)}

    % Unit normal vector u
    \PacketInkPath[->,black!65,line width=0.55pt]{(0,0) -- (0,1) node[left=3pt,font=\small,text=black!80] {\InkLabel{$u$}}}

    \coordinate (O) at (0,0);
    \coordinate (W) at (2.40,0);
    \coordinate (V) at (2.40,1.42);
    \coordinate (HV) at (2.40,-1.42);
    \coordinate (UperpEnd) at (2.95,0);
    \coordinate (SpanUEnd) at (0,-1.40);

    % Focal vectors v (black), Hv (single restrained accent), projection w (mid-gray)
    \PacketInkPath[->,line width=0.55pt,black]{(O) -- (V)}
    \PacketInkPath[->,line width=0.55pt,figblue]{(O) -- (HV)}
    \PacketInkPath[->,line width=0.55pt,black!65,shorten >=1.8pt]{(O) -- (W)
      node[pos=0.58,sloped,below=2pt,font=\scriptsize,text=black!75] {\InkLabel{$w$}}}

    % Two equal normal steps -(u^T v)u (dotted)
    \draw[ink guide arrow,draw=black!60,shorten <=1.6pt,shorten >=2.2pt] (V) -- (W);
    \draw[ink guide arrow,draw=black!60,shorten <=2.2pt,shorten >=2.2pt] (W) -- (HV);

    \fill[black] (O) circle (1.3pt);
    \fill[black] (W) circle (1.3pt);
  \end{scope}

  % Unrotated crisp labels
  \node[right=4pt,font=\small] at (UperpEnd) {\InkLabel{$u^\perp$}};
  \node[below=2pt,font=\scriptsize,text=black!70] at (SpanUEnd) {\InkLabel{$\spn(u)$}};
  \node[above right=1pt,font=\small] at (V) {\InkLabel{$v$}};
  \node[below right=1pt,font=\small,text=figblue!90!black] at (HV) {\InkLabel{$Hv$}};
  \node[right=5pt,font=\scriptsize,text=black!75] at ($(V)!0.5!(W)$) {\InkLabel{$-(u^{\mathsf T}v)u$}};
  \node[right=5pt,font=\scriptsize,text=black!75] at ($(W)!0.5!(HV)$) {\InkLabel{$-(u^{\mathsf T}v)u$}};
\end{tikzpicture}
\figcap{Householder reflection across $u^\perp$: with $w=P_{u^\perp}v$ fixed (eigenvalue $+1$ on $u^\perp$), subtracting twice the normal component $(u^{\mathsf T}v)u$ along $\spn(u)$ (eigenvalue $-1$) sends $v$ to its mirror image $Hv$.}
\end{fig}

\begin{proposition*}[QR and Cholesky factorizations, LADR 7D]\target{name-qr}\target{name-cholesky}
  \begin{enumerate}
    \ii[(a)] \textbf{(QR)} If $A$ is $m\times n$ with linearly independent columns $v_1,\dots,v_n$, then $A=QR$ where $Q$ ($m\times n$) has orthonormal columns and $R$ ($n\times n$) is upper-triangular with positive diagonal entries.
    \ii[(b)] \textbf{(Cholesky)} Every positive definite $n\times n$ matrix $B$ factors uniquely as $B=R^*R$ where $R$ is upper-triangular with positive diagonal entries.
  \end{enumerate}
\end{proposition*}
\begin{proof}
  \begin{steps}
    \ii \textbf{(a)} Apply Gram--Schmidt \ilink{T2}{(T2)} to $v_1,\dots,v_n$ to obtain orthonormal $e_1,\dots,e_n$ with $\spn(e_1,\dots,e_k)=\spn(v_1,\dots,v_k)$. Let $Q=[e_1\;\cdots\;e_n]$, set $R_{jk}=\langle v_k,e_j\rangle$ ($j\le k$) and $0$ ($j>k$). Then column $k$ of $QR$ is $\sum_{j\le k}\langle v_k,e_j\rangle e_j=v_k$ \ilink{T2}{(T2)}, and with Gram--Schmidt orthogonal remainder $w_k$, $R_{kk}=\langle v_k,e_k\rangle=\|w_k\|>0$.
    \ii \textbf{(b)} By Prop.~106, $B=M^*M$ for an invertible $M$. Applying QR gives $M=QR$ with $Q$ unitary ($Q^*Q=I$) and $R$ upper-triangular with positive diagonal, whence $B=(QR)^*QR=R^*Q^*QR=R^*R$.
  \end{steps}
\end{proof}

\noindent\textbf{The question answered.} $S(x,y)=(x,y,0)$ has orthonormal columns, so it is an isometry (Prop.~110). Its adjoint $S^*(a,b,c)=(a,b)$ annihilates $(0,0,1)$, so $S^*$ is \emph{not} an isometry ($S^*S=I_2$ while $SS^*=\operatorname{diag}(1,1,0)\neq I_3$). On the same space $\RR^n$, yes: $S^*S=I$ forces $SS^*=I$ by the inverse rule.

\begin{shortcuts}
  \ii \textbf{Checking isometry:} Verify columns are orthonormal (Prop.~110).
  \ii \textbf{Inverse of isometric operator:} $S^{-1}=S^*$ with zero row reduction.
  \ii \textbf{Cholesky ($2\times2$):} For positive definite $B=\begin{psmallmatrix}b_{11}&b_{12}\\\overline{b_{12}}&b_{22}\end{psmallmatrix}$, $R=\begin{psmallmatrix}r_{11}&r_{12}\\0&r_{22}\end{psmallmatrix}$ has $r_{11}=\sqrt{b_{11}}$, $r_{12}=b_{12}/r_{11}$, $r_{22}=\sqrt{b_{22}-|r_{12}|^2}$.
\end{shortcuts}

\begin{trapbox}
  \ii \textbf{Orthogonal columns $\neq$ orthonormal:} $\operatorname{diag}(2,1)$ has orthogonal columns yet is not an isometry.
  \ii \textbf{$S^*S=I$ does not force $SS^*=I$ when $V\neq W$:} $S\colon\RR^2\to\RR^3$ above has $S^*S=I_2$ yet $SS^*\neq I_3$.
  \ii \textbf{Prop.~111 needs $\CC$:} Over $\RR$, $2\begin{psmallmatrix}0&-1\\1&0\end{psmallmatrix}$ has no real eigenvalues at all yet doubles every length.
\end{trapbox}

\subsection{Polar decomposition (Ch.~22)}

\begin{qbox}
  A nonzero complex scalar factors as $z=e^{i\theta}|z|$: a rotation $e^{i\theta}$ times a positive stretch $|z|=\sqrt{\overline{z}z}$. Does every linear operator $T\in\LL(V)$ factor similarly as an isometry times a positive operator, $T=S\sqrt{T^*T}$? Find this factorization for
  \[
    A=\begin{bmatrix}-1&-2\\2&1\end{bmatrix}.
  \]
\end{qbox}

Notice why $\sqrt{T^*T}$ is the natural ``modulus $|T|$'' of $T$: $T^*T$ is always positive ($\langle T^*Tv,v\rangle=\|Tv\|^2\ge0$), so $P=\sqrt{T^*T}$ exists uniquely (Prop.~107), and for every $v\in V$:
\begin{equation*}\tag{$\ast$}
  \|\sqrt{T^*T}\,v\|^2 = \langle\sqrt{T^*T}\,v,\sqrt{T^*T}\,v\rangle = \langle T^*Tv,v\rangle = \|Tv\|^2
  \implies \|\sqrt{T^*T}\,v\| = \|Tv\|.
\end{equation*}
Equation~$(\ast)$ says $Pv$ and $Tv$ always have the same length.

\begin{theorem*}[Prop.~112 / LADR 7.93: Polar decomposition]\target{prop-112}\target{name-polar}
  For every operator $T\in\LL(V)$, there exists an isometry $S\in\LL(V)$ such that
  \[
    T = S\sqrt{T^*T}.
  \]
  Moreover, $P=\sqrt{T^*T}$ is uniquely determined by $T=SP$ ($S$ isometry, $P$ positive); and if $T$ is invertible, $S=TP^{-1}$ is unique as well.
\end{theorem*}
\begin{ideabox}
  Because $\|Pv\|=\|Tv\|$ by $(\ast)$, the map sending $Pv\mapsto Tv$ is well-defined and preserves all lengths from $\range P$ onto $\range T$. By Rank--Nullity and \ilink{T3}{(T3)}, $(\range P)^\perp$ and $(\range T)^\perp$ have equal dimension, so we match them with any isometry between orthonormal bases.
\end{ideabox}
\begin{proof}
  Write $P=\sqrt{T^*T}$.
  \begin{steps}
    \ii Define $S_1\colon\range P\to\range T$ by $S_1(Pv)=Tv$. If $Pv_1=Pv_2$, then $\|T(v_1-v_2)\|=\|P(v_1-v_2)\|=0$ by $(\ast)$, so $Tv_1=Tv_2$; hence $S_1$ is well-defined and clearly linear, onto $\range T$, and length-preserving ($\|S_1(Pv)\|=\|Tv\|=\|Pv\|$). Thus $\dim\range P=\dim\range T$ \ilink{T4}{(T4)}.
    \ii Consequently $\dim(\range P)^\perp = \dim(\range T)^\perp = k$ \ilink{T3}{(T3)}. Choose orthonormal bases $g_1,\dots,g_k$ of $(\range P)^\perp$ and $h_1,\dots,h_k$ of $(\range T)^\perp$, and let $S_2\colon(\range P)^\perp\to(\range T)^\perp$ be the linear map $S_2g_j=h_j$ (isometric by T2).
    \ii Decomposing $v=u+w$ along $\range P\oplus(\range P)^\perp$ ($u=P_{\range P}v$, $w=v-u$), define $Sv=S_1u+S_2w$. Linearity holds because $P_{\range P}$ is linear \ilink{T3}{(T3)}; since $S_1u\in\range T\perp(\range T)^\perp\ni S_2w$, Pythagoras \ilink{T1}{(T1)} gives $\|Sv\|^2=\|S_1u\|^2+\|S_2w\|^2=\|u\|^2+\|w\|^2=\|v\|^2$, so $S$ is an isometry and $S(Pv)=S_1(Pv)=Tv$ for all $v$.
    \ii \emph{Uniqueness:} If $T=SP$ ($S$ isometry, $P$ positive), then $T^*T=P^*S^*SP=P^2$ (Prop.~109), forcing $P=\sqrt{T^*T}$ (Prop.~107). If $T$ is invertible, $P$ is invertible (by $(\ast)$, $Pv=0\implies Tv=0\implies v=0$) and $S=TP^{-1}$ is unique. (When $T$ is invertible, setting $S=TP^{-1}$ directly avoids Steps~1--3 since $S^*S=P^{-1}T^*TP^{-1}=P^{-1}P^2P^{-1}=I$.)
  \end{steps}
\end{proof}

\begin{example*}[Ex.~77: Polar form on $\CC$]\target{ex-77}
  Multiplication $z\mapsto\lambda z$ ($\lambda\neq0$) has adjoint $w\mapsto\overline{\lambda}w$, so $\sqrt{\lambda^*\lambda}=\sqrt{|\lambda|^2}=|\lambda|$ and $\lambda = (\lambda/|\lambda|)|\lambda| = e^{i\theta}|\lambda|$.
\end{example*}

\begin{example*}[Worked examples: invertible vs.\ non-invertible]
  \begin{itemize}
    \ii \emph{The question ($A=\begin{psmallmatrix}-1&-2\\2&1\end{psmallmatrix}$):} $A^{\mathsf T}A=\begin{psmallmatrix}5&4\\4&5\end{psmallmatrix}$ has $\tr=10$, $\det=9$. By the $2\times2$ square-root formula,
    \[
      P=\sqrt{A^{\mathsf T}A}=\frac{A^{\mathsf T}A+3I}{\sqrt{10+6}}=\begin{bmatrix}2&1\\1&2\end{bmatrix},\qquad
      S = AP^{-1} = \frac{1}{3}\begin{bmatrix}-1&-2\\2&1\end{bmatrix}\begin{bmatrix}2&-1\\-1&2\end{bmatrix} = \begin{bmatrix}0&-1\\1&0\end{bmatrix},
    \]
    a counterclockwise $90^\circ$ rotation. So $A$ stretches by $3$ along $(1,1)$ and $1$ along $(1,-1)$, then rotates by $90^\circ$.
    \ii \emph{Non-invertible $T=\begin{psmallmatrix}0&1\\0&0\end{psmallmatrix}$:} Here $T^{\mathsf T}T=\operatorname{diag}(0,1)$ gives $P=\operatorname{diag}(0,1)$, $\range P=\spn(e_2)$, $\range T=\spn(e_1)$, $S_1e_2=e_1$, while $S_2\colon\spn(e_1)\to\spn(e_2)$ can send $e_1\mapsto e_2$ \emph{or} $e_1\mapsto -e_2$:
    \[
      S=\begin{bmatrix}0&1\\\pm1&0\end{bmatrix}, \qquad SP=T.
    \]
    So $S$ is \emph{not} unique when $T$ is non-invertible.
  \end{itemize}
\end{example*}

\begin{fig}
\begin{tikzpicture}[scale=1.0]
  % Uniform scale=0.52 across all 3 panels
  \begin{scope}[xshift=1.4cm,yshift=1.7cm,scale=0.52]
  \InkRegion[inkblue]{(0,0) circle (1)}
    \PacketInkPath[figaxis,draw=black!35]{(-1.8,0) -- (1.9,0)}
    \PacketInkPath[figaxis,draw=black!35]{(0,-1.8) -- (0,1.9)}
      \PacketInkPath[black,line width=0.55pt]{(0,0) circle (1)}
    \PacketInkPath[black!55,line width=0.38pt,rotate=-45]{(0.34,0) -- (0.34,0.34) -- (0,0.34)}
    \PacketInkPath[->,line width=0.55pt,figblue]{(0,0) -- ({1/sqrt(2)},{1/sqrt(2)})
      node[above right=1pt,font=\scriptsize,text=figblue!90!black] {\InkLabel{$e_1$}}}
    \PacketInkPath[->,line width=0.55pt,black]{(0,0) -- ({1/sqrt(2)},{-1/sqrt(2)})
      node[below right=1pt,font=\scriptsize] {\InkLabel{$e_2$}}}
  \end{scope}
  \node[font=\footnotesize,anchor=north] at (1.4,0.12) {\InkLabel{(a)}};

  \PacketInkPath[->,line width=0.55pt,black!75]{(2.65,1.7) -- node[above=2pt,font=\footnotesize] {\InkLabel{$P=\sqrt{A^{\mathsf T}A}$}} (4.35,1.7)}

  \begin{scope}[xshift=6.4cm,yshift=1.7cm,scale=0.52]
  \InkRegion[inkblue]{(0,0) [rotate=45] ellipse (3 and 1)}
    \PacketInkPath[figaxis,draw=black!35]{(-2.9,0) -- (3.0,0)}
    \PacketInkPath[figaxis,draw=black!35]{(0,-2.6) -- (0,2.75)}
      \draw[ink guide,draw=black!45] (0,0) circle (1);
    \PacketInkPath[black,line width=0.55pt,rotate=45]{(0,0) ellipse (3 and 1)}
    \PacketInkPath[black!55,line width=0.38pt,rotate=-45]{(0.34,0) -- (0.34,0.34) -- (0,0.34)}
    \PacketInkPath[->,line width=0.55pt,figblue]{(0,0) -- ({3/sqrt(2)},{3/sqrt(2)})
      node[above right=1pt,font=\scriptsize,text=figblue!90!black] {\InkLabel{$3e_1$}}}
    \PacketInkPath[->,line width=0.55pt,black]{(0,0) -- ({1/sqrt(2)},{-1/sqrt(2)})
      node[below right=2pt,font=\scriptsize] {\InkLabel{$e_2$}}}
  \end{scope}
  \node[font=\footnotesize,anchor=north] at (6.4,0.12) {\InkLabel{(b)}};

  \PacketInkPath[->,line width=0.55pt,black!75]{(8.25,1.7) -- node[above=2pt,font=\footnotesize] {\InkLabel{$S$}} (9.95,1.7)}

  \begin{scope}[xshift=12.0cm,yshift=1.7cm,scale=0.52]
  \InkRegion[inkblue]{(0,0) [rotate=135] ellipse (3 and 1)}
    \PacketInkPath[figaxis,draw=black!35]{(-2.9,0) -- (3.0,0)}
    \PacketInkPath[figaxis,draw=black!35]{(0,-2.6) -- (0,2.75)}
      \draw[ink guide,draw=black!45] (0,0) circle (1);
    \PacketInkPath[black,line width=0.55pt,rotate=135]{(0,0) ellipse (3 and 1)}
    \PacketInkPath[black!55,line width=0.38pt,rotate=45]{(0.34,0) -- (0.34,0.34) -- (0,0.34)}
    \PacketInkPath[->,line width=0.55pt,figblue]{(0,0) -- ({-3/sqrt(2)},{3/sqrt(2)})
      node[above left=1pt,font=\scriptsize,text=figblue!90!black] {\InkLabel{$3f_1$}}}
    \PacketInkPath[->,line width=0.55pt,black]{(0,0) -- ({1/sqrt(2)},{1/sqrt(2)})
      node[above right=2pt,font=\scriptsize] {\InkLabel{$f_2$}}}
  \end{scope}
  \node[font=\footnotesize,anchor=north] at (12.0,0.12) {\InkLabel{(c)}};
\end{tikzpicture}
\figcap{Polar decomposition and SVD $A=SP$ for $A=\begin{psmallmatrix}-1&-2\\2&1\end{psmallmatrix}$ at uniform scale (dotted unit circle): $P=\sqrt{A^{\mathsf T}A}$ stretches along $e_1,e_2$ by $s_1=3,s_2=1$, and the rotation $S$ sends $e_j\mapsto f_j$.}
\end{fig}

\begin{proposition*}[Real $2\times2$ polar formula \& determinant product rule]\target{name-polar-2x2}
  For $2\times2$ matrices $X,Y$, $\det(XY)=\det X\det Y$. Consequently, if $A=\begin{psmallmatrix}a&b\\c&d\end{psmallmatrix}$ is real with $\det A\neq0$ and polar decomposition $A=SP$, then $P=S^{\mathsf T}A$ and:
  \begin{enumerate}
    \ii[(a)] if $\det A>0$, $S=\frac{1}{N_1}\begin{bmatrix}a+d&b-c\\c-b&a+d\end{bmatrix}$ with $N_1=\sqrt{(a+d)^2+(c-b)^2}$ (a rotation);
    \ii[(b)] if $\det A<0$, $S=\frac{1}{N_2}\begin{bmatrix}a-d&b+c\\b+c&d-a\end{bmatrix}$ with $N_2=\sqrt{(a-d)^2+(b+c)^2}$ (a reflection).
  \end{enumerate}
\end{proposition*}
\begin{proof}
  Expanding $(x_1y_1+x_2y_3)(x_3y_2+x_4y_4)-(x_1y_2+x_2y_4)(x_3y_1+x_4y_3)=(x_1x_4-x_2x_3)(y_1y_4-y_2y_3)$ proves $\det(XY)=\det X\det Y$. In (a), $\det A>0$ prevents $a+d=c-b=0$ (else $\det A=-a^2-b^2\le0$), $S$ has orthonormal columns with $\det S=1$, and $S^{\mathsf T}A$ has off-diagonal difference $\bigl((a+d)(b-c)+(c-b)(a+d)\bigr)/N_1=0$ (symmetric), trace $N_1>0$, and $\det P=\det S^{\mathsf T}\det A=\det A>0$, so $P$ is positive definite (Prop.~106) and $SP=SS^{\mathsf T}A=A$. Case (b) follows identically with $\det S=-1$ and trace $N_2>0$.
\end{proof}

\begin{trapbox}
  \ii \textbf{Right positive factor uses $T^*T$, not $TT^*$:} If you factor $T=P'S$ instead, $P'=\sqrt{TT^*}$.
  \ii \textbf{$\sqrt{T^*T}$ is not entrywise $|T|$:} For $A=\begin{psmallmatrix}-1&-2\\2&1\end{psmallmatrix}$, $P=\begin{psmallmatrix}2&1\\1&2\end{psmallmatrix}$, not $\begin{psmallmatrix}1&2\\2&1\end{psmallmatrix}$.
\end{trapbox}

\subsection{Singular values, SVD, and the pseudoinverse (Ch.~22, LADR 7E--7F)}

\begin{qbox}
  Returning to $A=\begin{psmallmatrix}-1&-2\\2&1\end{psmallmatrix}$, which unit vector $v\in\RR^2$ maximizes $\|Av\|$, and how large is the maximum? Furthermore, what is the closest rank-$1$ matrix to $A$?
\end{qbox}

For rectangular maps $T\in\LL(V,W)$, eigenvalues do not even make sense because $T$ maps between \emph{different} spaces. Just as least-squares replaces $Ax=b$ with $A^{\mathsf T}Ax=A^{\mathsf T}b$, we study the positive self-adjoint operator $T^*T\in\LL(V)$.

\begin{proposition*}[LADR 7.43 / 7E: Properties of $T^*T$]\target{name-ladr-7-43}
  For any $T\in\LL(V,W)$: $T^*T\in\LL(V)$ is positive, $\nul(T^*T)=\nul T$, $\range(T^*T)=\range T^*$, and $\dim\range(T^*T)=\dim\range T$.
\end{proposition*}
\begin{proof}
  Positivity is $\langle T^*Tv,v\rangle=\|Tv\|^2\ge0$. Clearly $Tv=0\implies T^*Tv=0$; conversely $T^*Tv=0\implies \|Tv\|^2=\langle T^*Tv,v\rangle=0\implies Tv=0$, so $\nul(T^*T)=\nul T$. Clearly $\range(T^*T)\subseteq\range T^*$, while Rank--Nullity \ilink{T4}{(T4)} and \ilink{T5}{(T5)} give $\dim\range(T^*T)=n-\dim\nul(T^*T)=n-\dim\nul T=\dim\range T=\dim\range T^*$, forcing equality.
\end{proof}

\begin{definition*}[Def.~54 / LADR 7.65: Singular values]\target{def-54}
  The \textbf{singular values} of $T\in\LL(V,W)$ are the eigenvalues of $\sqrt{T^*T}$ listed with multiplicity---equivalently, the nonnegative square roots of the eigenvalues of $T^*T$. There are $n=\dim V$ of them, ordered $s_1\ge s_2\ge\dots\ge s_n\ge0$.
\end{definition*}

\begin{proposition*}[Positive singular values, LADR 7E]
  For $T\in\LL(V,W)$, the number of strictly positive singular values equals $\dim\range T=\rank T$. Thus $T$ is injective $\iff 0$ is not a singular value of $T$.
\end{proposition*}
\begin{proof}
  In an orthonormal eigenbasis $e_1,\dots,e_n$ of $\sqrt{T^*T}$ with $\sqrt{T^*T}e_j=s_je_j$, $\range\sqrt{T^*T}=\spn\{e_j:s_j>0\}$ has dimension equal to $\#\{j:s_j>0\}$, and by $(\ast)$, $\nul\sqrt{T^*T}=\nul T$, so Rank--Nullity gives $\dim\range T=\dim\range\sqrt{T^*T}$.
\end{proof}

\begin{theorem*}[Prop.~113 / LADR 7.70 \& 7.69: Singular value decomposition]\target{prop-113}\target{name-svd}\target{name-ladr-7-69}
  Let $T\in\LL(V)$ have singular values $s_1\ge\dots\ge s_n$. There exist orthonormal bases $e_1,\dots,e_n$ and $f_1,\dots,f_n$ of $V$ such that $Te_j=s_jf_j$, hence
  \[
    Tv = \sum_{j=1}^n s_j\langle v,e_j\rangle f_j \qquad\text{and}\qquad T^*w = \sum_{j=1}^n s_j\langle w,f_j\rangle e_j \qquad \text{(LADR 7.69)}
  \]
  for all $v,w\in V$. In matrix form, $A = F\Sigma E^* = \sum_{j=1}^n s_j f_j e_j^*$ with $E,F$ unitary and $\Sigma=\operatorname{diag}(s_1,\dots,s_n)$.
\end{theorem*}
\begin{proof}
  Choose an orthonormal eigenbasis $e_1,\dots,e_n$ of $\sqrt{T^*T}$ with $\sqrt{T^*T}e_j=s_j e_j$ \ilink{T7}{(T7)}, and take $T=S\sqrt{T^*T}$ (Prop.~112). Set $f_j=Se_j$; by Prop.~110, $f_1,\dots,f_n$ is an orthonormal basis of $V$, and $Te_j=S(s_j e_j)=s_j f_j$. Expanding $v=\sum_j\langle v,e_j\rangle e_j$ \ilink{T2}{(T2)} gives $Tv=\sum_j s_j\langle v,e_j\rangle f_j$. For $T^*w$, testing against $e_k$ gives $\langle T^*w,e_k\rangle = \langle w,Te_k\rangle = \langle w,s_k f_k\rangle = s_k\langle w,f_k\rangle$, so Fourier expansion in $\{e_k\}$ yields LADR~7.69: $T^*w = \sum_k s_k\langle w,f_k\rangle e_k$ (and $T^*f_k=s_ke_k$).
\end{proof}

\begin{proposition*}[Prop.~116: Converse---diagonal in two ONBs means singular values]\target{prop-116}
  If $e_1,\dots,e_n$ and $f_1,\dots,f_n$ are orthonormal bases of $V$ with $Te_k=s_k f_k$ and $s_k\ge0$ for all $k$, then $s_1,\dots,s_n$ are the singular values of $T$.
\end{proposition*}
\begin{proof}
  For all $j,k$, $\langle e_j,T^*f_k\rangle=\langle Te_j,f_k\rangle=s_j\delta_{jk}=\langle e_j,s_k e_k\rangle$, so $T^*f_k=s_k e_k$ \ilink{T5}{(T5)}, hence $T^*Te_k=s_k^2 e_k$. Defining $Re_k=s_k e_k$ gives a positive operator (Prop.~105) with $R^2=T^*T$, so $R=\sqrt{T^*T}$ (Prop.~107) has eigenvalues $s_1,\dots,s_n$ with multiplicity \ilink{T8}{(T8)}.
\end{proof}

\begin{proposition*}[Rectangular SVD between two spaces, $T\in\LL(V,W)$]
  Let $T\in\LL(V,W)$ have singular values $s_1\ge\dots\ge s_r>0=s_{r+1}=\dots=s_n$ ($r=\rank T$). There exist orthonormal bases $e_1,\dots,e_n$ of $V$ and $f_1,\dots,f_m$ of $W$ such that $Te_j=s_j f_j$ ($j\le r$) and $Te_j=0$ ($j>r$), hence
  \[
    Tv = \sum_{j=1}^r s_j\langle v,e_j\rangle f_j, \qquad
    T^*w = \sum_{j=1}^r s_j\langle w,f_j\rangle e_j, \qquad
    A = F\Sigma E^*,
  \]
  where $\Sigma$ is $m\times n$ with $s_1,\dots,s_r$ on its leading diagonal and $0$ elsewhere.
\end{proposition*}
\begin{proof}
  Diagonalize positive $T^*T$ on $V$ by an orthonormal eigenbasis $e_1,\dots,e_n$ with $T^*Te_j=s_j^2 e_j$. Then $\langle Te_i,Te_j\rangle=s_i^2\delta_{ij}$, so $f_j=Te_j/s_j$ ($j\le r$) are orthonormal in $W$, while $\|Te_j\|^2=0\implies Te_j=0$ ($j>r$). Extend $f_1,\dots,f_r$ to an orthonormal basis $f_1,\dots,f_m$ of $W$ \ilink{T2}{(T2)}. Checking $\langle T^*w,e_j\rangle=\langle w,Te_j\rangle$ yields $T^*w=\sum_{j=1}^r s_j\langle w,f_j\rangle e_j$.
\end{proof}

\begin{example*}[Three concrete SVDs + Ex.~78]\target{ex-78}
  \begin{enumerate}
    \ii \emph{Zero singular value:} $A=\begin{psmallmatrix}2&2\\1&1\end{psmallmatrix}$ gives $A^{\mathsf T}A=\begin{psmallmatrix}5&5\\5&5\end{psmallmatrix}$, eigenvalues $10,0$, $s_1=\sqrt{10},s_2=0$; $e_1=\frac{1}{\sqrt{2}}(1,1)$, $e_2=\frac{1}{\sqrt{2}}(1,-1)$, $f_1=\frac{Ae_1}{\sqrt{10}}=\frac{1}{\sqrt{5}}(2,1)$, $f_2=\frac{1}{\sqrt{5}}(-1,2)$, so $A=\sqrt{10}\,f_1 e_1^{\mathsf T}$.
    \ii \emph{$3\times2$ matrix:} for
    \[
      A = \begin{bmatrix}1&1\\0&1\\1&0\end{bmatrix},
    \]
    we get $A^{\mathsf T}A=\begin{psmallmatrix}2&1\\1&2\end{psmallmatrix}$ ($s_1=\sqrt{3},s_2=1$, $e_1=\frac{1}{\sqrt{2}}(1,1)$, $e_2=\frac{1}{\sqrt{2}}(1,-1)$). Then $f_1=\frac{1}{\sqrt{6}}(2,1,1)$, $f_2=\frac{1}{\sqrt{2}}(0,-1,1)$, completed by $f_3=\frac{1}{\sqrt{3}}(1,-1,-1)$.
    \ii \emph{Ex.~78 (Notes Ch.~22):} $A=\begin{psmallmatrix}1&0\\1&1\end{psmallmatrix}$ has $A^{\mathsf T}A=\begin{psmallmatrix}2&1\\1&1\end{psmallmatrix}$ with $\tr=3$, $\det=1$, eigenvalues $\frac{3\pm\sqrt{5}}{2}=\varphi^{\pm2}$ where $\varphi=\frac{1+\sqrt{5}}{2}$. Thus $s_1=\varphi\approx1.618$, $s_2=\varphi^{-1}=\varphi-1\approx0.618$ (even though the only eigenvalue of $A$ is $1$), $c=\sqrt{\varphi^2+1}$, $e_1=\frac{1}{c}(\varphi,1)$, $e_2=\frac{1}{c}(-1,\varphi)$, $f_1=\frac{1}{c}(1,\varphi)$, $f_2=\frac{1}{c}(-\varphi,1)$, and polar rotation $S=\frac{1}{\sqrt{5}}\begin{psmallmatrix}2&-1\\1&2\end{psmallmatrix}$.
  \end{enumerate}
\end{example*}

\noindent\textbf{The spectral norm, smallest stretch, and ellipsoid image (LADR 7F).} The largest singular value $s_1$ is the \textbf{spectral norm} $\|T\|_2$.

\begin{theorem*}[Prop.~114 \& LADR 7F: Largest/smallest stretch, products, and ellipsoid image]\target{prop-114}\target{name-ladr-7-75}
  Let $S,T\in\LL(V)$ with singular values of $T$ ordered $s_1\ge\dots\ge s_n$.
  \begin{enumerate}
    \ii[(a)] $\|Tv\|\le s_1\|v\|$ for all $v\in V$, with equality at $v=e_1$, so $\|T\|_2 = \max_{\|v\|=1}\|Tv\| = s_1$.
    \ii[(b)] On $\LL(V)$, $\|\cdot\|_2$ is a norm (positivity, definiteness, scaling, triangle inequality), yet if $\dim V\ge2$, it does \textbf{not} come from any inner product.
    \ii[(c)] $\|Tv\|\ge s_n\|v\|$ for all $v\in V$, with equality at $v=e_n$; $\|ST\|_2\le\|S\|_2\|T\|_2$.
    \ii[(d)] If $T$ is invertible, the singular values of $T^{-1}$ are $1/s_n\ge\dots\ge 1/s_1$ ($\|T^{-1}\|_2=1/s_n$), and $T$ maps the unit sphere $\{\|v\|=1\}$ onto the ellipsoid $\bigl\{\sum_{j=1}^n d_j f_j : \sum_{j=1}^n |d_j|^2/s_j^2 = 1\bigr\}$ with orthogonal semi-axes $s_1f_1,\dots,s_nf_n$ (\textbf{LADR 7.75}).
  \end{enumerate}
\end{theorem*}
\begin{proof}
  \begin{steps}
    \ii \textbf{(a),(c).} Writing $v=\sum_j c_j e_j$ in Prop.~113, $Tv=\sum_j s_j c_j f_j$ so $s_n^2\sum|c_j|^2 \le \|Tv\|^2 = \sum s_j^2|c_j|^2 \le s_1^2\sum|c_j|^2 = s_1^2\|v\|^2$, with equality at $v=e_1$ and $v=e_n$. Also $\|ST\|_2 = \max_{\|v\|=1}\|S(Tv)\| \le \|S\|_2\max_{\|v\|=1}\|Tv\| = \|S\|_2\|T\|_2$.
    \ii \textbf{(b).} Definiteness: $\|T\|_2=0\implies\|Tv\|\le0\ \forall v\implies T=0$; homogeneity and triangle inequality follow from max over $\|v\|=1$. For $\dim V\ge2$, $E_{11}=\operatorname{diag}(1,0,\dots)$ and $E_{22}=\operatorname{diag}(0,1,0,\dots)$ satisfy $\|E_{11}\|_2=\|E_{22}\|_2=\|E_{11}+E_{22}\|_2=\|E_{11}-E_{22}\|_2=1$ (Prop.~116), violating the parallelogram law $1^2+1^2\neq 2(1^2+1^2)$ \ilink{T10}{(T10)}.
    \ii \textbf{(d).} From $Te_k=s_kf_k$ ($s_k>0$), $T^{-1}f_k=s_k^{-1}e_k$, so Prop.~116 gives singular values $s_k^{-1}$. And $v=\sum c_je_j$ has $\|v\|^2=\sum|c_j|^2=1 \iff Tv=\sum d_jf_j$ with $d_j=s_jc_j$ satisfies $\sum|d_j|^2/s_j^2=1$.
  \end{steps}
\end{proof}

\begin{theorem*}[Prop.~115: Eckart--Young best rank-$k$ approximation]\target{prop-115}\target{name-eckart-young}
  Let $T\in\LL(V)$ have singular values $s_1\ge\dots\ge s_n$ and SVD $Tv=\sum_{j=1}^n s_j\langle v,e_j\rangle f_j$. For $1\le k<n$, define the rank-$\le k$ truncation
  \[
    T_k v = \sum_{j=1}^k s_j\langle v,e_j\rangle f_j.
  \]
  Then $\rank T_k\le k$, $\|T-T_k\|_2 = s_{k+1}$, and for every $R\in\LL(V)$ with $\rank R\le k$, $\|T-R\|_2 \ge s_{k+1}$. Thus $T_k$ is a best rank-$\le k$ approximation to $T$ in spectral norm.
\end{theorem*}
\begin{ideabox}
  The error $T-T_k$ drops the first $k$ modes and keeps $s_{k+1},\dots,s_n$, so its largest stretch is $s_{k+1}$. No rank-$\le k$ map $R$ can achieve smaller error: in the $(k+1)$-dimensional subspace $U=\spn(e_1,\dots,e_{k+1})$, Rank--Nullity prevents $R|_U$ from being injective, so some unit vector $x\in U$ lies in $\nul R$; on that vector $x$, $(T-R)x = Tx$ has norm at least $s_{k+1}$.
\end{ideabox}
\begin{proof}
  \begin{steps}
    \ii $\range T_k\subseteq\spn(f_1,\dots,f_k)$, so $\rank T_k\le k$.
    \ii $(T-T_k)v = \sum_{j>k}s_j c_j f_j$ has $\|(T-T_k)v\|^2=\sum_{j>k}s_j^2|c_j|^2\le s_{k+1}^2\|v\|^2$, with equality at $e_{k+1}$, so $\|T-T_k\|_2=s_{k+1}$.
    \ii For any $R$ with $\rank R\le k$, set $U=\spn(e_1,\dots,e_{k+1})$. Since $\dim U=k+1 > k \ge \dim\range R$, Rank--Nullity \ilink{T4}{(T4)} yields a unit vector $x=\sum_{j=1}^{k+1}c_je_j\in U\cap\nul R$ ($\sum_{j=1}^{k+1}|c_j|^2=1$). Then $Rx=0$ gives
    \[
      \|T-R\|_2 \ge \|(T-R)x\| = \|Tx\| = \left(\sum_{j=1}^{k+1} s_j^2 |c_j|^2\right)^{1/2} \ge s_{k+1}.
    \]
  \end{steps}
\end{proof}

\begin{fig}
\begin{tikzpicture}[scale=1.0]
  \begin{axis}[
    at={(0,0)},anchor=south west,
    ybar,bar width=16pt,bar shift=0pt,
    width=9.2cm,height=4.7cm,
    ymin=0,ymax=4.85,
    xmin=0.4,xmax=5.6,
    xtick={1,2,3,4,5},
    xticklabels={\InkLabel{$s_1$},\InkLabel{$s_2$},\InkLabel{$s_3$},\InkLabel{$s_4$},\InkLabel{$s_5$}},
    ytick={0,2,4},
    extra y ticks={0.95},
    extra y tick labels={\color{figred!90!black}\InkLabel{$s_3$}},
    extra y tick style={tick style={draw=figred!85!black}},
    axis line style={draw=black!75,line width=0.45pt},
    tick style={draw=black!60,line width=0.38pt},
    tick label style={font=\scriptsize},
    label style={font=\footnotesize},
    ylabel={\InkLabel{$s_j$}},
    ylabel style={rotate=-90},
    axis lines*=left,
    clip=false
  ]
    % Vertical cutoff at k = 2
    \draw[ink guide,draw=inkslate] (axis cs:2.5,0) -- (axis cs:2.5,4.65);

    % Kept modes (soft blue wash), leading discarded mode s_3 (single accent), tail (white hairline)
    \addplot[pattern=inkgrain,pattern color=inkblue,draw=inkblue,line width=.45pt] coordinates {(1,4.0) (2,2.4)};
    \addplot[pattern=inkgrain,pattern color=inkred,draw=inkred,line width=.45pt] coordinates {(3,0.95)};
    \addplot[ink plot,fill=white,draw=black!65,line width=0.45pt] coordinates {(4,0.45) (5,0.2)};

    % Spectral-norm error threshold line at s_3
    \draw[ink guide,draw=inkred] (axis cs:0.4,0.95) -- (axis cs:5.55,0.95);

    % Quiet atomic group labels
    \node[font=\scriptsize,text=black!75] at (axis cs:1.50,4.55) {\InkLabel{$j\le k$}};
    \node[font=\scriptsize,text=black!65] at (axis cs:4.05,4.55) {\InkLabel{$j>k$}};

    % Clean error label seated above threshold line
    \node[font=\scriptsize,text=figred!90!black,anchor=south] at (axis cs:4.50,1.08) {\InkLabel{$\|T-T_k\|_2 = s_3$}};
  \end{axis}
\end{tikzpicture}
\figcap{Eckart--Young best rank-$k$ approximation ($k=2$): keeping $s_1,s_2$ (blue) leaves the discarded tail $j>k$, whose largest value $s_3$ sets the exact spectral-norm error $\|T-T_k\|_2=s_3$.}
\end{fig}

\begin{theorem*}[Pseudoinverse $T^\dagger$ via SVD \& Least-Squares, LADR 7.99, 7.101, 7.111]\target{name-pseudoinverse}
  Let $T\in\LL(V,W)$ have nonzero singular values $s_1,\dots,s_r>0$ and SVD $Tv=\sum_{j=1}^r s_j\langle v,e_j\rangle f_j$. Define the \textbf{pseudoinverse} $T^\dagger\in\LL(W,V)$ by
  \[
    T^\dagger w = \sum_{j=1}^r \frac{1}{s_j}\langle w,f_j\rangle e_j \qquad (\text{matrix form: } A^\dagger = E\Sigma^\dagger F^*,\; \Sigma^\dagger_{jj}=s_j^{-1}\ (j\le r),\; 0\text{ else}).
  \]
  Then:
  \begin{enumerate}
    \ii[(a)] $T^\dagger T = P_{(\nul T)^\perp}$ (orthogonal projection onto $(\nul T)^\perp=\range T^*$) and $TT^\dagger = P_{\range T}$;
    \ii[(b)] For any $b\in W$, $x^+ = T^\dagger b$ is a least-squares solution minimizing $\|Tx-b\|$, and among all least-squares minimizers $x^+$ is the \textbf{unique} one of smallest norm $\|x\|$.
  \end{enumerate}
\end{theorem*}
\begin{proof}
  Note $(\nul T)^\perp = \spn(e_1,\dots,e_r)$ and $\range T = \spn(f_1,\dots,f_r)$. Directly from the formulas: $T^\dagger T e_j = e_j$ ($j\le r$) and $0$ ($j>r$), proving $T^\dagger T=P_{(\nul T)^\perp}$; similarly $TT^\dagger f_j = f_j$ ($j\le r$) and $0$ ($j>r$), proving $TT^\dagger=P_{\range T}$. For $b\in W$, $Tx^+=TT^\dagger b=P_{\range T}b$ is the closest point to $b$ in $\range T$ (\xlink{math110}{2026-11-02-orthogonality-projections}{name-best-approximation}{Week~10}). Any other minimizer has the form $x=x^+ + z$ with $z\in\nul T$. Since $x^+\in(\nul T)^\perp$, Pythagoras gives $\|x\|^2=\|x^+\|^2+\|z\|^2\ge\|x^+\|^2$ with strict inequality unless $z=0$.
\end{proof}

\noindent\textbf{The question answered.} The maximum stretch over unit vectors is $s_1=3$ at $v=e_1=\frac{1}{\sqrt{2}}(1,1)$ (Prop.~114). The closest rank-$1$ matrix in spectral norm is $A^{(1)}=3f_1e_1^{\mathsf T}=\frac{3}{2}\begin{psmallmatrix}-1&-1\\1&1\end{psmallmatrix}$, with error $\|A-A^{(1)}\|_2=s_2=1$ (Prop.~115).

\begin{shortcuts}
  \ii \textbf{Sum of squares:} $\sum_{j=1}^n s_j^2 = \tr(A^*A) = \sum_{i,j}|a_{ij}|^2$ \ilink{T11}{(T11)}.
  \ii \textbf{$2\times2$ singular values:} $s_1^2+s_2^2=\sum|a_{ij}|^2$ and $s_1s_2=|\det A|$, determining $s_1,s_2$ directly from two numbers.
  \ii \textbf{Rank-$1$ map ($A=uv^*$):} Singular values are $\|u\|\,\|v\|,0,\dots,0$ with $e_1=v/\|v\|$ and $f_1=u/\|u\|$.
  \ii \textbf{Use the smaller of $A^*A$ and $AA^*$:} Both have the same nonzero eigenvalues $s_j^2$. If $m<n$, find $f_j$ from $AA^*$ and set $e_j=A^*f_j/s_j$.
  \ii \textbf{Singular values of $T^*$, $ST$, $TS$:} Identical to those of $T$ whenever $S$ is an isometric operator.
\end{shortcuts}

\begin{trapbox}
  \ii \textbf{Singular values $\neq |\lambda_j|$ unless normal:} Ex.~78 has eigenvalues $1,1$ yet singular values $\varphi\neq\varphi^{-1}$.
  \ii \textbf{Compute $f_j=Ae_j/s_j$ directly:} Independently normalizing eigenvectors of $AA^*$ risks sign flips ($-f_1$ gives $Ae_1=-s_1f_1$).
  \ii \textbf{Singular values are $\sqrt{\lambda_j(T^*T)}$, not $\lambda_j(T^*T)$:} And for $m\times n$ matrices, $\Sigma$ is $m\times n$ (not square).
\end{trapbox}

\section{Problem set}

\begin{prob}\target{prob-HW12-P1}
  \textbf{HW12 P1.}\ Classify each real symmetric matrix as positive definite, positive semidefinite but not positive definite, or not positive semidefinite:
  \[
    A=\begin{bmatrix}2&1\\1&2\end{bmatrix},\qquad
    B=\begin{bmatrix}1&2\\2&1\end{bmatrix},\qquad
    C=\begin{bmatrix}1&1\\1&1\end{bmatrix}.
  \]

\end{prob}

\begin{prob}\target{prob-HW12-P2}
  \textbf{HW12 P2.}\ Find the positive square root of $A=\begin{bmatrix}2&1\\1&2\end{bmatrix}$ as an explicit $2\times2$ matrix.

\end{prob}

\begin{prob}\target{prob-HW12-P3}
  \textbf{HW12 P3.}\ Let $S\colon V\to W$ be linear between finite-dimensional inner product spaces, and $e_1,\dots,e_n$ an orthonormal basis of $V$. Prove that $S$ is an isometry if and only if $Se_1,\dots,Se_n$ is an orthonormal list in $W$.

\end{prob}

\begin{prob}\target{prob-HW12-P4}
  \textbf{HW12 P4.}\ Find the polar decomposition $A=S\sqrt{A^*A}$ of $A=\begin{bmatrix}0&2\\1&0\end{bmatrix}$ acting on $\RR^2$ with the standard inner product.

\end{prob}

\begin{prob}\target{prob-HW12-P5}
  \textbf{HW12 P5.}\ Find a singular value decomposition of $A=\begin{bmatrix}1&1\\1&1\end{bmatrix}$: give the singular values and orthonormal bases $e_1,e_2$ and $f_1,f_2$ with $Ae_j=s_jf_j$.

\end{prob}

\begin{prob}\target{prob-HW12-P6}
  \textbf{HW12 P6 (optional).}\ Let $T$ be an invertible operator. Prove that the singular values of $T^{-1}$ are the reciprocals of the singular values of $T$.

\end{prob}

\begin{prob}\target{prob-HW12-P7}
  \textbf{HW12 P7 (optional).}\ Let $T$ have singular values $s_1\ge\dots\ge s_n\ge0$. Prove that $\|Tv\|\le s_1\|v\|$ for every $v$, and that equality holds for some $v\neq0$.

\end{prob}

\begin{prob}\target{prob-O1}
  \textbf{O1 (commuting positive operators).}\ Let $S,T\in\LL(V)$ be positive with $ST=TS$. Prove that $ST$ is positive. Show by an example that $ST$ need not be positive if $ST\neq TS$.

\end{prob}

\begin{prob}\target{prob-O2}
  \textbf{O2 (a $2\times3$ matrix).}\ Find a singular value decomposition $A=F\Sigma E^{\mathsf T}$ of $A=\begin{bmatrix}1&0&1\\0&1&1\end{bmatrix}$.

\end{prob}

\begin{prob}\target{prob-O3}
  \textbf{O3 (polar decomposition and SVD).}\ Find the polar decomposition and a singular value decomposition of $A=\begin{bmatrix}3&0\\4&5\end{bmatrix}$.

\end{prob}

\begin{prob}\target{prob-O4}
  \textbf{O4 (norm of $T^*T$).}\ Prove that $\|T^*T\|_2=\|T\|_2^2$.

\end{prob}

\clearpage
\section*{Hints}
Try each problem before consulting its hint.

\par\medskip\noindent\textbf{\hyperlink{prob-HW12-P1}{HW12 P1}} \ilink{prop-105}{Props.~105--106}; $\det$ and $\tr$ are the product and sum of the eigenvalues \ilink{T9}{(T9)}.

\par\medskip\noindent\textbf{\hyperlink{prob-HW12-P2}{HW12 P2}} The \ilink{name-square-root-2x2}{square-root formula}; note $4+2\sqrt{3}=(1+\sqrt{3})^2$.

\par\medskip\noindent\textbf{\hyperlink{prob-HW12-P3}{HW12 P3}} \ilink{prop-110}{Prop.~110}: (i)$\Rightarrow$(ii)$\Rightarrow$(iii) and (iv)$\Rightarrow$(i).

\par\medskip\noindent\textbf{\hyperlink{prob-HW12-P4}{HW12 P4}} $A^{\mathsf T}A$ is diagonal; then $S=AP^{-1}$.

\par\medskip\noindent\textbf{\hyperlink{prob-HW12-P5}{HW12 P5}} Follow the SVD method; one singular value is $0$, so extend $\{f_1\}$ to an orthonormal basis.

\par\medskip\noindent\textbf{\hyperlink{prob-HW12-P6}{HW12 P6 (optional)}} From $Te_k=s_kf_k$, solve for $T^{-1}f_k$; then \ilink{prop-116}{Prop.~116}.

\par\medskip\noindent\textbf{\hyperlink{prob-HW12-P7}{HW12 P7 (optional)}} \ilink{prop-113}{Prop.~113} and the calculation in \ilink{prop-114}{Prop.~114}.

\par\medskip\noindent\textbf{\hyperlink{prob-O1}{O1 (commuting positive operators)}} $T$ preserves each eigenspace of $S$, on which $\sqrt{S}$ is a scalar multiple of $I$ (\ilink{prop-107}{Prop.~107}); hence $\sqrt{S}$ commutes with $T$, and $\langle STv,v\rangle=\langle T\sqrt{S}v,\sqrt{S}v\rangle$.

\par\medskip\noindent\textbf{\hyperlink{prob-O2}{O2 (a $2\times3$ matrix)}} $AA^{\mathsf T}$ is only $2\times2$; find its unit eigenvectors $f_j$ first, then set $e_j=A^{\mathsf T}f_j/s_j$.

\par\medskip\noindent\textbf{\hyperlink{prob-O3}{O3 (polar decomposition and SVD)}} $A^{\mathsf T}A$ has $\tr=50$ and $\det=225$; apply the \ilink{name-square-root-2x2}{square-root formula} or \ilink{name-polar-2x2}{polar formula}.

\par\medskip\noindent\textbf{\hyperlink{prob-O4}{O4 (norm of $T^*T$)}} $T^*T$ is positive, so its singular values coincide with its eigenvalues $s_j^2$.

\end{document}
